\part{The Typewriter}

\chapter{The Fixed Character Array}

The typewriter replaced continuous handwriting with selection from a fixed set. Once writing is selection, it can be taught as a table of positions.

\section{From Writing Ball to Sholes}

Several mechanical writing devices preceded the commercially successful typewriter. The machine developed in Milwaukee by Christopher Latham Sholes with Carlos Glidden and Samuel Soule was patented in 1868 and manufactured under E.~Remington and Sons from 1873. \needs{verify patent year and manufacturing start against a museum or patent source} Type bars struck the paper from below, so the operator could not see what had just been typed, a condition that shaped early training: the learner had to trust position and not appearance.

\section{The Consolidation of QWERTY}

The QWERTY arrangement stabilized in the 1870s and became the default through manufacturer competition and the growth of trained typists. The popular explanation that the order was designed to prevent adjacent type bars from jamming is widely repeated and has been disputed in the literature. \needs{cite the dispute; do not adjudicate here} For the purposes of this book what matters is that a large population trained on one arrangement created a cost to changing it, a cost that Chapter~12 examines through the Dvorak case.

\section{The Shift Key}

The addition of a shift mechanism allowed one key to produce two characters. The consequence for pedagogy is that capitals are a separate skill: the learner must coordinate a second hand with the first. The binder examined in Chapter~8 introduces capitals by alternating case within a drill line (\texttt{AsDsWs}), which isolates exactly this coordination.

\chapter{Touch Typing and the Commercial School}

Touch typing is the practice of typing without looking at the keys, using all fingers according to a fixed assignment. It is not a property of the typewriter. It is a training doctrine.

\section{The Ten-Finger Method}

Accounts of the origin of touch typing usually mention the memorization of the keyboard by court stenographers and the 1888 speed contest in Cincinnati between Frank McGurrin and Louis Traub, which was reported as a victory for the touch method. \needs{primary or high-quality secondary account of the 1888 contest and of earlier claimants to the method} Whatever the precise history, the doctrine that emerged prescribes a home position, an assignment of every key to one finger, and a habit of not looking.

\section{Home Row as Doctrine}

Home-row doctrine defines a resting position for the fingers and treats every other key as a reach from it. This is why the first lesson of a large family of tutors is the same eight keys, and why the reach drills of later lessons are written as excursions from and returns to home. Chapter~8 shows this structure in two independent printed sources.

\section{The Commercial School}

The vocational school gave the doctrine an institutional form: a fixed curriculum, a supervised room, a machine of known type, and an examiner. The apparatus of Chapter~8 is a late instance of this form. \needs{history of commercial schools and business colleges in the period}

\chapter{Office Labor and Ergonomic Doctrine}

Typing became an occupation, and an occupation acquires rules for the body.

\section{The Typing Pool and the Contest}

Large offices concentrated typists in pools, and typing contests and trade demonstrations publicized speed as a marketable skill. \needs{sources on typing pools, contests, and their gendered labor history} Speed figures functioned as credentials, which is why so many tutors report their results in words per minute even when accuracy would predict employability better.

\section{Posture and Rhythm}

Instructional literature prescribed posture, wrist position, and sitting height, and it also prescribed rhythm. Typing to music and to recorded pacing signals was promoted as a way of regularizing keystroke intervals. \needs{documentation of rhythm-typing records and their dates} The idea of an external pacing signal is a through-line of this book: the phonograph record, the metronome, the cassette tape of Chapter~8, the metronome of Chapter~19, and the falling letters of Chapter~20 are successive answers to one problem, namely how to make a learner keep a rate that the learner does not yet control.

\section{Discipline}

The typing lesson is a discipline in the ordinary sense: a routine that shapes a body by repetition under observation. The paper tutor makes this explicit, because its supervisor is the observer and its typed page is the record.

\chapter{The Paper Tutor}
\label{ch:paper}

This chapter reconstructs, from photographed pages, the rules of two printed typing curricula. Neither contains a processor. Each is nonetheless a program in the sense that matters here: an ordered sequence of instructions, a fixed data format, and a runtime made of people and machines.

\section{Two Printed Sources}

Two distinct sources are used. They must be kept apart.

\begin{description}
\item[The Academy pages.] Pages bearing the imprint of Academy of Learning, among them pages 1, 11, 12, 20, 21, 31 and 34, with ``Revised 07/92'' printed on page 34. The pages are organized as numbered \emph{Exercises}.
\item[The workbook.] A set titled \emph{Keyboard Skill Building}, copyright Qwerty Learning Systems Inc.\ 1992, code KBD103/392, organized as twenty \emph{Lessons} with a table of contents, instructions to students, progress tests, and cassette tape practice.
\end{description}

Both carry 1992 dates and both are typewriter-based, so they may have been used in the same institution. \inferred{the shared date and medium suggest a common institutional setting; the pages do not say so} The book does not treat them as one document. Whether they shared a physical binder, and how a learner moved between them, is unresolved and is recorded in Appendix~F.

\section{The Academy Pages}

\subsection{The apparatus}

Page 1 is a labeled diagram of an electronic daisy-wheel typewriter, naming the platen knobs, paper bail, paper release lever, ribbon cassette, daisy wheel cassette, printing position indicator, and carrying handle. Its caption states that the learner's typewriter should look the same or very similar to the one illustrated, and that otherwise the learner should ask the supervisor for the manual for the machine in use.

Three consequences follow. Hardware is part of the specification. Variation in hardware is delegated to a human. And the pedagogy resides in the pages, while everything that adapts to circumstance resides in the supervisor.

\subsection{The page specification}

Each exercise page opens with a keyboard diagram in which the new keys are filled black and earlier keys remain outlined, a line naming the new keys and their fingers, a left margin of 1 to 1.5 inches, and single spacing. Exercise~1 assigns A, S, D, F to the left hand, semicolon, L, K, J to the right hand, and the space bar to the right thumb. The header fixes the physical output as well as the input: the learner is producing a correctly formatted page, which is the vocational deliverable, while acquiring the keystrokes.

\subsection{The key order}

The order in which keys are introduced is not alphabetical and not by letter frequency. On the pages available it is geometric.

\begin{center}
\begin{tabular}{@{}clp{6.5cm}@{}}
\toprule
Exercise & New keys & Position \\
\midrule
1 & A S D F ; L K J, space & home row, both hands \\
2 & G H & home row, index stretch \\
6 & W I & upper row, first reaches \\
13 & Z & lower row, left little finger \\
\bottomrule
\end{tabular}
\end{center}

Exercises 3 to 5 and 7 to 12 are not among the pages available. The home row is completed before any reach is attempted, and the rare letter Z arrives at Exercise~13, so its position follows finger geometry, not usefulness.

\subsection{The early groups}

Each exercise page is divided into numbered groups, each printed as three identical lines. On the pages for Exercises 1, 2, 6 and 13 the first six groups follow a common pattern, transcribed in Appendix~C.

\begin{enumerate}[label=\textbf{G\arabic*.}]
\item A chain of the new keys, for example \texttt{asdf ;lkj} (Exercise~1) and \texttt{asdfg ;lkjh} (Exercise~2). The right hand is written in mirror order.
\item The chain with one old key inserted, for example \texttt{asdfa ;lkj;} and \texttt{asdfgfa ;lkjhj;}. From Exercise~6 this group also carries the first shift practice: \texttt{AsDsWs ;LkIkL}. Exercise~13 gives \texttt{aSaZa ;L;P;}.
\item Either a scramble of the first group (\texttt{dafs k;jl}, Exercise~1) or short words (\texttt{wed kid was lie}, Exercise~6; \texttt{zip zoo zig zag}, Exercise~13).
\item Word list. \texttt{ass lad add alf}; \texttt{glad asks hash gala gald}; \texttt{ail law his owe}; \texttt{zinc: zulu zero jazz:}.
\item Longer words. \texttt{flak alas fall asks}; \texttt{salad glass shall flash jaffa slash}; \texttt{were soil will hide}; \texttt{quiz doze size hazy:}.
\item A phrase or a list of longer words. \texttt{a sad lad asks dad;} closes Exercise~2, and \texttt{heir wear like wore} and \texttt{zeal lazy fizz fuzz:} close Exercises 6 and 13.
\end{enumerate}

Group~3 of Exercise~1 contains one deliberate variant (\texttt{k;jk} among \texttt{k;jl}), so the line is not a perfect repetition.

\subsection{The later groups}

Two pages show that exercises continue past Group~6, and that the later groups are sentences.

Page 21 continues Exercise~6 with Groups 7 to 11. Group 7 is a word chain (\texttt{wise lied wild wire lie}), Group 8 is a short sentence (\texttt{A wood fire is good; oil is worse for a fire; wood is for fire}), and Groups 9 to 11 are longer sentences in which each content word draws on the new keys:

\begin{Verbatim}
Group 9    We killed wild doe for a while for food; Will would ride well
Group 10   A fellow said he walked a high hill; he grew riled if we rode;
Group 11   Ilse who is a redhead is a selfish girl who likes a wild life;
\end{Verbatim}

Page 31 belongs to an exercise between Exercises 6 and 13 whose number is not established from the photograph. Its Groups 7 to 13 run from word chains to sentences:

\begin{Verbatim}
Group 9    Peter and Paul ran quickly to the printers for a paper:
Group 10   King Pepe snapped his fingers as the Pekinese pup passed by;
Group 12   I will quit the room again if he does not answer all queries:
Group 13   The Queen was presented with two quills at the banquet today;
\end{Verbatim}

So the group count varies by exercise, and the six-group ladder is the early portion of a longer sequence that ends in continuous sentences. The last group of each observed exercise is a sentence or a run of sentences, and the sentences are selected to exercise the new keys. The page for Exercise~13 and other pages show faint lines through the paper from the reverse side, and these are not quoted as text.

\section{The Sandwich Drill}

Both sources teach a reach by the same device: the reached key is placed between two presses of the home-row key that normally covers it. In the Academy pages the device appears as \texttt{asdsws ;lkikl} (Exercise~6, Group~1). In the workbook it opens the second block of Lesson~1:

\begin{Verbatim}
aqa sws ded frf ftf fgf ;p; lol kik juj jyj jhj aza sxs
dcd fvf fbf ;/; 1.1 k,k jmj jnj
\end{Verbatim}

and is immediately followed by a line pairing each drill with a word that contains it (\texttt{aqa aqua}, \texttt{sws sail}, \texttt{ded dead}, \texttt{frf rain}, \texttt{ftf tall}, \texttt{fgf gulf}). The sandwich drill isolates the transition from home to reach and back, which is the unit that the learner state of Chapter~17 calls transition competence.

\section{A Formal Model of the Academy Pages}

Let $\Kset_i$ be the set of keys admitted by exercise $i$, with
\[
\Kset_0\subset \Kset_1\subset\cdots\subset \Kset_n,
\]
and let $\Newk_i=\Kset_i\setminus \Kset_{i-1}$ be the new keys. For a word $w$ let $\Sig(w)$ be the set of characters it requires. The vocabulary available to the author of exercise $i$ is
\[
\Vocab_i=\{\,w\in\Vocab:\Sig(w)\subseteq \Kset_i\,\}.
\]
Two regimes appear. In the earliest exercises the binding constraint is admissibility itself: $\Vocab_1$ is so small that the lesson is nearly all of it, which is why \texttt{alf} and \texttt{gald} appear. In later exercises $\Vocab_i$ is large and the author selects for salience,
\[
\Vocab_i^{\mathrm{new}}=\{\,w\in\Vocab_i:\Sig(w)\cap\Newk_i\neq\varnothing\,\},
\]
which is why Exercise~13 is a catalogue of words containing Z and why every word in Exercise~6 contains W or I. The constraint migrates from what can be spelled to what will exercise the new key.

Each group $G_j$ has a cyclic unit $u_j$ and is printed as three lines, each a window of width $\omega$ on the endless repetition of the unit:
\[
\mathrm{line}_j=\mathrm{prefix}_{\omega}\bigl(u_j^{\infty}\bigr),\qquad G_j=(\mathrm{line}_j)^{3}.
\]
Lines end mid-unit (\texttt{asd}, \texttt{daf}, \texttt{lag}), so the carriage return never coincides with a unit boundary and the typist has no natural place to pause. An exercise is the sequence $E_i=(G_1,\dots,G_m)$ with $m$ varying by exercise. What is not on the pages is the advancement rule, and the judgment $J$ that moves the learner from $E_i$ to $E_{i+1}$ is not written in the Academy pages available.

\section{The Workbook}

\subsection{Lesson template}

The workbook uses a different template. Its lessons do not restrict the alphabet: Lesson~1 opens with a speed drill of ordinary full-alphabet sentences. Each lesson is a fixed sequence of sections with a stated number of repetitions.

\begin{center}
\begin{tabular}{@{}lc p{7.2cm}@{}}
\toprule
Section & Repeats & Content in Lesson 1 \\
\midrule
Speed Drill & 3 & four ordinary sentences, each typed three times \\
Keyboard Drill & 2 & sandwich drills, then drills paired with words \\
Alphabet Drill & 2 & alphabet in runs, then one three-letter word per letter, then a full-alphabet sentence \\
Word families, phrases, sentence & 1 & inflected word families, slash-separated phrases, four sentences \\
Tape practice & once & paced typing at a stated rate \\
Timed writings & 2 & five-minute tests \\
\bottomrule
\end{tabular}
\end{center}

The repetition counts fall (3, 2, 2, 1) as the material becomes longer and more varied, a schedule different from the uniform three lines of the Academy pages.

\subsection{The lesson schedule}

The table of contents gives the sequence of the twenty lessons. Every fifth lesson is a test lesson with no tape practice or timed writings; the others end with tape practice at a rate that steps upward in blocks.

\begin{center}
\begin{tabular}{@{}lll@{}}
\toprule
Lessons & Drill emphasis & Tape rate \\
\midrule
1 to 4 & keyboard; word families, prefixes, suffixes & 25 wpm \\
5 & practice for and administration of Progress Test 1 & none \\
6 to 9 & number drill; suffixes, double letters & 30 wpm \\
10 & Progress Test 2 & none \\
11 to 14 & symbol drill; word families, prefixes, suffixes, double letters & 35 wpm \\
15 & Progress Test 3 & none \\
16 to 19 & alphabet sentences; alternate-hand, one-hand words, spelling & 40 wpm \\
20 & Final Progress Test & none \\
\bottomrule
\end{tabular}
\end{center}

The workbook therefore has a numeric speed curriculum: a stepped schedule of five words per minute per block, delivered by a cassette tape acting as a pacer. This is absent from the Academy pages.

\subsection{Two policies on errors}

The workbook gives its errors policy in printed instructions. The general instructions say that typing errors are \emph{not} to be corrected. Drill lines are then marked and counted: the sample page shows errors circled by the student and a count written beside each set of lines, and a stated scale reads zero to three errors as excellent, four to seven as satisfactory, and eight or more as a signal to practice further. The phrase lines contain diagonal slash marks which the learner must not type; the instructions say they are present to make the learner think and type in groups of words.

The cassette practice reverses the priority. Its instructions tell the learner not to worry about the number of errors and to concentrate on keeping up with the tape, so that speed is developed. The one workbook therefore contains two theories at once: accuracy is the object of the drill lines and pace is the object of the tape. Chapter~15 develops this.

Because errors are not corrected, the page records every uncorrected error by policy, not only by the limits of the machine. This restores, for the workbook, the claim that a typed page is a complete record of uncorrected error. It does not record timing, and it cannot record corrections, because none are permitted.

\subsection{Timed writings, progress tests, and the completion sheet}

Timed writings are five minutes long. The learner practices each passage once, then is timed for five minutes, repeated twice. Numbers at the right of the passage give the cumulative word count for each line, and a second set of numbers beneath the timing is used when the test ends mid-line; neither set is to be typed. Paragraphs are indented five spaces and line endings must match the passage so that the test can be scored. The learner uses a clean sheet, types a name on the sixth line from the top, triple spaces, sets double spacing, and types one timing on the front of the sheet and the second on the back.

The printed text on the back of these instruction pages shows through the paper and is only partly legible. It appears to say that the learner may do each test twice, hands in the better one to be checked and initialed, and, when it is returned, enters the speed on the Lesson Completion Sheet and files the test as the last sheet of that set of lessons. \inferred{this reading is from show-through and must be checked against the reverse page before quotation}

\section{The Judgment $J$}

In the Academy pages the advancement rule $J$ is unwritten, and its absence is evidence that it lives in a person. In the workbook $J$ is partly documented. A learner completes lessons in sequence; scores drill lines against a stated scale; produces two five-minute timings per lesson; has the better one checked and initialed by a supervisor; enters a speed on a completion sheet; and is tested at Lessons 5, 10, 15 and 20. What remains unknown is whether a threshold on the completion sheet gated advancement. The Lesson Completion Sheet and a Progress Test page are needed to settle this. \needs{Lesson Completion Sheet and one Progress Test page}

\section{The Paper Tutor as a Distributed System}

Assembled, the workbook and the Academy pages together have components that map onto the parts of a program.

\begin{center}
\begin{tabular}{@{}p{4.2cm}p{7.4cm}@{}}
\toprule
Program component & Paper tutor \\
\midrule
Instruction sequence & ordered exercises or lessons \\
Data format & header block, groups or sections, repeated lines \\
Execution engine & learner's hands and the typewriter \\
Pacing signal & cassette tape (workbook only) \\
Output record & typed page, uncorrected \\
Error counter & the learner, marking and counting \\
Conditional branch & supervisor's check, initials, and test results \\
Exception handler & supervisor, for machine variation \\
Persistent state & the filed stack of lessons and the completion sheet \\
\bottomrule
\end{tabular}
\end{center}

The program is executed across several substrates at once. Its state after any session is the stack of typed pages, which records every uncorrected error but no intervals between keystrokes.

\section{What Passed to the DOS Tutor}

Automated tutors preserved the graded key order, the per-lesson lexicon, the highlighted keyboard diagram, and short-unit repetition, and they absorbed the pacing tape into the computer. They added measured speed, immediate error feedback, and stored scores. They generally discarded the page. A DOS tutor commonly keeps a number where the paper tutor kept an artifact, and the practice ledger of Chapter~27 is a deliberate return to the page.

\section*{Figures}
\begin{itemize}
\item Figure 8.1: typewriter diagram and supervisor caption (Academy p.~1).
\item Figure 8.2: Exercise 1, home row (p.~11).
\item Figure 8.3: Exercise 2, G and H (p.~12).
\item Figure 8.4: Exercise 6, W and I, with first shift alternation (pp.~20, 21).
\item Figure 8.5: Exercise 13, Z, with the ``Revised 07/92'' imprint (p.~34).
\item Figure 8.6: a later exercise ending in sentences (p.~31).
\item Figure 8.7: workbook instruction pages (drill lines, cassette practice, timed writings).
\item Figure 8.8: workbook table of contents, Lessons 1 to 20.
\item Figure 8.9: workbook Lesson~1, page 1-1.
\end{itemize}

\chapter*{Study Questions, Part II}
\addcontentsline{toc}{chapter}{Study Questions, Part II}
\begin{enumerate}
\item Explain why a machine on which the typist cannot see the line just typed favors training by position over training by appearance.
\item State the sandwich drill as a rule. Why does it isolate the transition between home and reach?
\item Using the transcriptions in Appendix~C, compute $\Vocab_i^{\mathrm{new}}$ for Exercise~6 and state which words on page~21 are outside it.
\item The workbook forbids correcting errors. List two things this policy makes observable and two that it hides.
\item Draw the supervisor's decision procedure for the workbook as far as the pages document it, and mark each edge whose rule is not printed.
\end{enumerate}
